\documentclass[10pt,a4paper]{amsart}
\usepackage{amsmath,amssymb,amsthm,graphicx}
\usepackage{verbatim}
\newenvironment{reference}{\comment}{\endcomment}
%\newenvironment{reference}{}{}
\newenvironment{slogan}{\comment}{\endcomment}
\newenvironment{history}{\comment}{\endcomment}
% For commutative diagrams we use Xy-pic
\usepackage[all]{xy}

% We use 2cell for 2-commutative diagrams.
\xyoption{2cell}
% Theorem environments.
%
\theoremstyle{plain}
\newtheorem{theorem}[subsection]{Theorem}
\newtheorem{proposition}[subsection]{Proposition}
\newtheorem{lemma}[subsection]{Lemma}

\theoremstyle{definition}
\newtheorem{definition}[subsection]{Definition}
\newtheorem{example}[subsection]{Example}
\newtheorem{exercise}[subsection]{Exercise}
\newtheorem{situation}[subsection]{Situation}

\theoremstyle{remark}
\newtheorem{remark}[subsection]{Remark}
\newtheorem{remarks}[subsection]{Remarks}

\numberwithin{equation}{subsection}

% Macros
%
\def\lim{\mathop{\mathrm{lim}}\nolimits}
\def\colim{\mathop{\mathrm{colim}}\nolimits}
\def\Spec{\mathop{\mathrm{Spec}}}
\def\Hom{\mathop{\mathrm{Hom}}\nolimits}
\def\Ext{\mathop{\mathrm{Ext}}\nolimits}
\def\SheafHom{\mathop{\mathcal{H}\!\mathit{om}}\nolimits}
\def\SheafExt{\mathop{\mathcal{E}\!\mathit{xt}}\nolimits}
\def\Sch{\mathit{Sch}}
\def\Mor{\mathop{\mathrm{Mor}}\nolimits}
\def\Ob{\mathop{\mathrm{Ob}}\nolimits}
\def\Sh{\mathop{\mathit{Sh}}\nolimits}
\def\NL{\mathop{N\!L}\nolimits}
\def\CH{\mathop{\mathrm{CH}}\nolimits}
\def\proetale{{pro\text{-}\acute{e}tale}}
\def\etale{{\acute{e}tale}}
\def\QCoh{\mathit{QCoh}}
\def\Ker{\mathop{\mathrm{Ker}}}
\def\Im{\mathop{\mathrm{Im}}}
\def\Coker{\mathop{\mathrm{Coker}}}
\def\Coim{\mathop{\mathrm{Coim}}}

% Boxtimes
%
\DeclareMathSymbol{\boxtimes}{\mathbin}{AMSa}{"02}

%
% Macros for moduli stacks/spaces
%
\def\QCohstack{\mathcal{QC}\!\mathit{oh}}
\def\Cohstack{\mathcal{C}\!\mathit{oh}}
\def\Spacesstack{\mathcal{S}\!\mathit{paces}}
\def\Quotfunctor{\mathrm{Quot}}
\def\Hilbfunctor{\mathrm{Hilb}}
\def\Curvesstack{\mathcal{C}\!\mathit{urves}}
\def\Polarizedstack{\mathcal{P}\!\mathit{olarized}}
\def\Complexesstack{\mathcal{C}\!\mathit{omplexes}}
% \Pic is the operator that assigns to X its picard group, usage \Pic(X)
% \Picardstack_{X/B} denotes the Picard stack of X over B
% \Picardfunctor_{X/B} denotes the Picard functor of X over B
\def\Pic{\mathop{\mathrm{Pic}}\nolimits}
\def\Picardstack{\mathcal{P}\!\mathit{ic}}
\def\Picardfunctor{\mathrm{Pic}}
\def\Deformationcategory{\mathcal{D}\!\mathit{ef}}
\usepackage{cleveref}
\crefname{originallemma}{original lemma}{original lemmas}
\Crefname{originallemma}{Original lemma}{Original lemmas}
\makeatletter
\newlabel{lemma-properties-diagonal-representable}{{lemma-properties-diagonal-representable}{source}{}{sourcepointer-1}{}}
\newlabel{lemma-properties-diagonal-representable@cref}{{[originallemma][0][]lemma-properties-diagonal-representable}{[1][0][]source}{}{}{}}
\makeatother
\crefname{originallemma}{original lemma}{original lemmas}
\Crefname{originallemma}{Original lemma}{Original lemmas}
\makeatletter
\newlabel{lemma-second-diagonal}{{lemma-second-diagonal}{source}{}{sourcepointer-2}{}}
\newlabel{lemma-second-diagonal@cref}{{[originallemma][0][]lemma-second-diagonal}{[1][0][]source}{}{}{}}
\makeatother
\crefname{originalreference}{original reference}{original references}
\Crefname{originalreference}{Original reference}{Original references}
\makeatletter
\newlabel{stacks-properties-lemma-monomorphism}{{stacks-properties-lemma-monomorphism}{source}{}{sourcepointer-3}{}}
\newlabel{stacks-properties-lemma-monomorphism@cref}{{[originalreference][0][]stacks-properties-lemma-monomorphism}{[1][0][]source}{}{}{}}
\makeatother
\title{Monomorphisms and representable morphisms are read off the first diagonal, and the second diagonal is always a monomorphism}
\author{The Stacks Project Authors}
\date{stacks-project:stacks-morphisms.tex:local-snapshot-20261012; GNU FDL 1.2 or later}
\begin{document}
\maketitle
\noindent\textit{Editorial note.} All delivered mathematical text is the author's own, copied from the pinned Stacks Project chapter source sources/stacks-algebraic-stacks-morphisms/pinned-original-source.tex; no statement or proof step was rewritten, reordered or removed.\par\medskip
\noindent\textit{Editorial note.} Preamble provenance: the Stacks Project keeps one single shared preamble for every chapter; this card takes it directly from the staged sources/stacks-algebraic-stacks-morphisms/preamble.tex as three verbatim spans (lines 8-12, 14-18 and 154-217), so no derived file was written and nothing was added or redefined. The wrapper class and the base packages are supplied by the assembler.\par\medskip
\noindent\textit{Editorial note.} Bibliography provenance: neither the counted statement nor its proof contains a citation (this lane rescans the whole span before writing the card), so the card carries no bibliography and no entry was reconstructed.\par\medskip
\noindent\textit{Editorial note.} Selection provenance: this card was opened under the file-based control contract (CONTROL.md), with the mechanical pre-card gate run on its statement and proof interval before the card was written.\par\medskip
\noindent\textit{Allowed standard prerequisites (editorial).}\begin{itemize}
\item The statement that a morphism of algebraic stacks is a monomorphism exactly when its diagonal is an isomorphism. Invoked at the start of the counted proof (source lines 1753-1755) as a cross-chapter pointer. Gives the first equivalence.
\item The statement that the second diagonal of a morphism of algebraic stacks is a monomorphism. Used inside the counted proof (source lines 1756-1758). Reformulates the second assertion as an if and only if for representability.
\item The characterization of representability of a morphism by its diagonal being a monomorphism. Invoked at the end of the counted proof (source lines 1759-1763). Turns the reformulation into the asserted equivalence.
\end{itemize}\medskip
\ifdefined\hypertarget\hypertarget{sourcepointer-1}{}\fi
\noindent\textit{Source reference.} Original lemma lemma-properties-diagonal-representable: sources/stacks-algebraic-stacks-morphisms/pinned-original-source.tex, source lines 519--519. Navigation pointer only: the counted text refers to the Stacks item lemma-properties-diagonal-representable, which belongs to another chapter of the project and is not part of this card. The pointer carries the target's own label as its number because this chapter's staged snapshot does not contain the other chapter's tag table; its content is not reproduced here.\par\medskip
\ifdefined\hypertarget\hypertarget{sourcepointer-2}{}\fi
\noindent\textit{Source reference.} Original lemma lemma-second-diagonal: sources/stacks-algebraic-stacks-morphisms/pinned-original-source.tex, source lines 1757--1757. Navigation pointer only: the counted text refers to the Stacks item lemma-second-diagonal, which belongs to another chapter of the project and is not part of this card. The pointer carries the target's own label as its number because this chapter's staged snapshot does not contain the other chapter's tag table; its content is not reproduced here.\par\medskip
\ifdefined\hypertarget\hypertarget{sourcepointer-3}{}\fi
\noindent\textit{Source reference.} Original reference stacks-properties-lemma-monomorphism: sources/stacks-algebraic-stacks-morphisms/pinned-original-source.tex, source lines 1167--1167. Navigation pointer only: the counted text refers to the Stacks item stacks-properties-lemma-monomorphism, which belongs to another chapter of the project and is not part of this card. The pointer carries the target's own label as its number because this chapter's staged snapshot does not contain the other chapter's tag table; its content is not reproduced here.\par\medskip
\begin{lemma}
\label{lemma-hierarchy}
A morphism $f : \mathcal{X} \to \mathcal{Y}$ of algebraic stacks is
\begin{enumerate}
\item a monomorphism if and only if $\Delta_{f, 1}$ is an isomorphism, and
\item representable by algebraic spaces if and only if $\Delta_{f, 1}$
is a monomorphism.
\end{enumerate}
Moreover, the second diagonal $\Delta_{f, 2}$ is always a monomorphism.
\end{lemma}

\begin{proof}
Recall from Properties of Stacks, Lemma
\ref{stacks-properties-lemma-monomorphism}
that a morphism of algebraic stacks is a monomorphism
if and only if its diagonal is an isomorphism of stacks.
Thus Lemma \ref{lemma-second-diagonal}
can be rephrased as saying that a morphism is
representable by algebraic spaces if the diagonal
is a monomorphism. In particular, it shows that condition
(3) of Lemma \ref{lemma-properties-diagonal-representable}
is actually an if and only if, i.e., a morphism of algebraic stacks
is representable by algebraic spaces if and only if
its diagonal is a monomorphism.
\end{proof}

\end{document}
